
\documentclass[12pt,a4paper,fleqn]{report}
\usepackage{latexsym,graphicx}
\usepackage[refpage]{nomencl}
\usepackage{../../iman,../../extra,../../isar,../../proof}
\usepackage{../../isabelle,../../isabellesym}
\usepackage{style}
\usepackage{../../pdfsetup}


\hyphenation{Isabelle}
\hyphenation{Isar}
\isadroptag{theory}

\title{\includegraphics[scale=0.5]{isabelle_isar}
  \\[4ex] Haskell-style type classes with Isabelle/Isar}
\author{\emph{Florian Haftmann}}

\begin{document}

\maketitle

\begin{abstract}
  This tutorial introduces the look-and-feel of Isar type classes
  to the end-user; Isar type classes are a convenient mechanism
  for organizing specifications, overcoming some drawbacks
  of raw axiomatic type classes. Essentially, they combine
  an operational aspect (in the manner of Haskell) with
  a logical aspect, both managed uniformly.
\end{abstract}

\thispagestyle{empty}\clearpage

\pagenumbering{roman}
\clearfirst

%
\begin{isabellebody}%
\def\isabellecontext{Classes}%
%
\isadelimtheory
%
\endisadelimtheory
%
\isatagtheory
\isacommand{theory}\isamarkupfalse%
\ Classes\isanewline
\isakeyword{imports}\ Main\ Setup\isanewline
\isakeyword{begin}%
\endisatagtheory
{\isafoldtheory}%
%
\isadelimtheory
%
\endisadelimtheory
%
\isamarkupsection{Introduction%
}
\isamarkuptrue%
%
\begin{isamarkuptext}%
Type classes were introduced by Wadler and Blott \cite{wadler89how}
  into the Haskell language to allow for a reasonable implementation
  of overloading\footnote{throughout this tutorial, we are referring
  to classical Haskell 1.0 type classes, not considering later
  additions in expressiveness}.  As a canonical example, a polymorphic
  equality function \isa{eq\ {\isasymColon}\ {\isasymalpha}\ {\isasymRightarrow}\ {\isasymalpha}\ {\isasymRightarrow}\ bool} which is overloaded on
  different types for \isa{{\isasymalpha}}, which is achieved by splitting
  introduction of the \isa{eq} function from its overloaded
  definitions by means of \isa{class} and \isa{instance}
  declarations: \footnote{syntax here is a kind of isabellized
  Haskell}

  \begin{quote}

  \noindent\isa{class\ eq\ where} \\
  \hspace*{2ex}\isa{eq\ {\isasymColon}\ {\isasymalpha}\ {\isasymRightarrow}\ {\isasymalpha}\ {\isasymRightarrow}\ bool}

  \medskip\noindent\isa{instance\ nat\ {\isasymColon}\ eq\ where} \\
  \hspace*{2ex}\isa{eq\ {\isadigit{0}}\ {\isadigit{0}}\ {\isacharequal}\ True} \\
  \hspace*{2ex}\isa{eq\ {\isadigit{0}}\ {\isacharunderscore}\ {\isacharequal}\ False} \\
  \hspace*{2ex}\isa{eq\ {\isacharunderscore}\ {\isadigit{0}}\ {\isacharequal}\ False} \\
  \hspace*{2ex}\isa{eq\ {\isacharparenleft}Suc\ n{\isacharparenright}\ {\isacharparenleft}Suc\ m{\isacharparenright}\ {\isacharequal}\ eq\ n\ m}

  \medskip\noindent\isa{instance\ {\isacharparenleft}{\isasymalpha}{\isasymColon}eq{\isacharcomma}\ {\isasymbeta}{\isasymColon}eq{\isacharparenright}\ pair\ {\isasymColon}\ eq\ where} \\
  \hspace*{2ex}\isa{eq\ {\isacharparenleft}x{\isadigit{1}}{\isacharcomma}\ y{\isadigit{1}}{\isacharparenright}\ {\isacharparenleft}x{\isadigit{2}}{\isacharcomma}\ y{\isadigit{2}}{\isacharparenright}\ {\isacharequal}\ eq\ x{\isadigit{1}}\ x{\isadigit{2}}\ {\isasymand}\ eq\ y{\isadigit{1}}\ y{\isadigit{2}}}

  \medskip\noindent\isa{class\ ord\ extends\ eq\ where} \\
  \hspace*{2ex}\isa{less{\isacharunderscore}eq\ {\isasymColon}\ {\isasymalpha}\ {\isasymRightarrow}\ {\isasymalpha}\ {\isasymRightarrow}\ bool} \\
  \hspace*{2ex}\isa{less\ {\isasymColon}\ {\isasymalpha}\ {\isasymRightarrow}\ {\isasymalpha}\ {\isasymRightarrow}\ bool}

  \end{quote}

  \noindent Type variables are annotated with (finitely many) classes;
  these annotations are assertions that a particular polymorphic type
  provides definitions for overloaded functions.

  Indeed, type classes not only allow for simple overloading but form
  a generic calculus, an instance of order-sorted algebra
  \cite{nipkow-sorts93,Nipkow-Prehofer:1993,Wenzel:1997:TPHOL}.

  From a software engineering point of view, type classes roughly
  correspond to interfaces in object-oriented languages like Java; so,
  it is naturally desirable that type classes do not only provide
  functions (class parameters) but also state specifications
  implementations must obey.  For example, the \isa{class\ eq}
  above could be given the following specification, demanding that
  \isa{class\ eq} is an equivalence relation obeying reflexivity,
  symmetry and transitivity:

  \begin{quote}

  \noindent\isa{class\ eq\ where} \\
  \hspace*{2ex}\isa{eq\ {\isasymColon}\ {\isasymalpha}\ {\isasymRightarrow}\ {\isasymalpha}\ {\isasymRightarrow}\ bool} \\
  \isa{satisfying} \\
  \hspace*{2ex}\isa{refl{\isacharcolon}\ eq\ x\ x} \\
  \hspace*{2ex}\isa{sym{\isacharcolon}\ eq\ x\ y\ {\isasymlongleftrightarrow}\ eq\ x\ y} \\
  \hspace*{2ex}\isa{trans{\isacharcolon}\ eq\ x\ y\ {\isasymand}\ eq\ y\ z\ {\isasymlongrightarrow}\ eq\ x\ z}

  \end{quote}

  \noindent From a theoretical point of view, type classes are
  lightweight modules; Haskell type classes may be emulated by SML
  functors \cite{classes_modules}.  Isabelle/Isar offers a discipline
  of type classes which brings all those aspects together:

  \begin{enumerate}
    \item specifying abstract parameters together with
       corresponding specifications,
    \item instantiating those abstract parameters by a particular
       type
    \item in connection with a ``less ad-hoc'' approach to overloading,
    \item with a direct link to the Isabelle module system:
      locales \cite{kammueller-locales}.
  \end{enumerate}

  \noindent Isar type classes also directly support code generation in
  a Haskell like fashion. Internally, they are mapped to more
  primitive Isabelle concepts \cite{Haftmann-Wenzel:2006:classes}.

  This tutorial demonstrates common elements of structured
  specifications and abstract reasoning with type classes by the
  algebraic hierarchy of semigroups, monoids and groups.  Our
  background theory is that of Isabelle/HOL \cite{isa-tutorial}, for
  which some familiarity is assumed.%
\end{isamarkuptext}%
\isamarkuptrue%
%
\isamarkupsection{A simple algebra example \label{sec:example}%
}
\isamarkuptrue%
%
\isamarkupsubsection{Class definition%
}
\isamarkuptrue%
%
\begin{isamarkuptext}%
Depending on an arbitrary type \isa{{\isasymalpha}}, class \isa{semigroup} introduces a binary operator \isa{{\isacharparenleft}{\isasymotimes}{\isacharparenright}} that is
  assumed to be associative:%
\end{isamarkuptext}%
\isamarkuptrue%
%
\isadelimquote
%
\endisadelimquote
%
\isatagquote
\isacommand{class}\isamarkupfalse%
\ semigroup\ {\isacharequal}\isanewline
\ \ \isakeyword{fixes}\ mult\ {\isacharcolon}{\isacharcolon}\ {\isachardoublequoteopen}{\isasymalpha}\ {\isasymRightarrow}\ {\isasymalpha}\ {\isasymRightarrow}\ {\isasymalpha}{\isachardoublequoteclose}\ \ \ \ {\isacharparenleft}\isakeyword{infixl}\ {\isachardoublequoteopen}{\isasymotimes}{\isachardoublequoteclose}\ {\isadigit{7}}{\isadigit{0}}{\isacharparenright}\isanewline
\ \ \isakeyword{assumes}\ assoc{\isacharcolon}\ {\isachardoublequoteopen}{\isacharparenleft}x\ {\isasymotimes}\ y{\isacharparenright}\ {\isasymotimes}\ z\ {\isacharequal}\ x\ {\isasymotimes}\ {\isacharparenleft}y\ {\isasymotimes}\ z{\isacharparenright}{\isachardoublequoteclose}%
\endisatagquote
{\isafoldquote}%
%
\isadelimquote
%
\endisadelimquote
%
\begin{isamarkuptext}%
\noindent This \hyperlink{command.class}{\mbox{\isa{\isacommand{class}}}} specification consists of two parts:
  the \qn{operational} part names the class parameter (\hyperlink{element.fixes}{\mbox{\isa{\isakeyword{fixes}}}}), the \qn{logical} part specifies properties on them
  (\hyperlink{element.assumes}{\mbox{\isa{\isakeyword{assumes}}}}).  The local \hyperlink{element.fixes}{\mbox{\isa{\isakeyword{fixes}}}} and \hyperlink{element.assumes}{\mbox{\isa{\isakeyword{assumes}}}} are lifted to the theory toplevel, yielding the global
  parameter \isa{{\isachardoublequote}mult\ {\isasymColon}\ {\isasymalpha}{\isasymColon}semigroup\ {\isasymRightarrow}\ {\isasymalpha}\ {\isasymRightarrow}\ {\isasymalpha}{\isachardoublequote}} and the
  global theorem \hyperlink{fact.semigroup.assoc:}{\mbox{\isa{semigroup{\isachardot}assoc{\isacharcolon}}}}~\isa{{\isachardoublequote}{\isasymAnd}x\ y\ z\ {\isasymColon}\ {\isasymalpha}{\isasymColon}semigroup{\isachardot}\ {\isacharparenleft}x\ {\isasymotimes}\ y{\isacharparenright}\ {\isasymotimes}\ z\ {\isacharequal}\ x\ {\isasymotimes}\ {\isacharparenleft}y\ {\isasymotimes}\ z{\isacharparenright}{\isachardoublequote}}.%
\end{isamarkuptext}%
\isamarkuptrue%
%
\isamarkupsubsection{Class instantiation \label{sec:class_inst}%
}
\isamarkuptrue%
%
\begin{isamarkuptext}%
The concrete type \isa{int} is made a \isa{semigroup} instance
  by providing a suitable definition for the class parameter \isa{{\isacharparenleft}{\isasymotimes}{\isacharparenright}} and a proof for the specification of \hyperlink{fact.assoc}{\mbox{\isa{assoc}}}.  This is
  accomplished by the \hyperlink{command.instantiation}{\mbox{\isa{\isacommand{instantiation}}}} target:%
\end{isamarkuptext}%
\isamarkuptrue%
%
\isadelimquote
%
\endisadelimquote
%
\isatagquote
\isacommand{instantiation}\isamarkupfalse%
\ int\ {\isacharcolon}{\isacharcolon}\ semigroup\isanewline
\isakeyword{begin}\isanewline
\isanewline
\isacommand{definition}\isamarkupfalse%
\isanewline
\ \ mult{\isacharunderscore}int{\isacharunderscore}def{\isacharcolon}\ {\isachardoublequoteopen}i\ {\isasymotimes}\ j\ {\isacharequal}\ i\ {\isacharplus}\ {\isacharparenleft}j{\isasymColon}int{\isacharparenright}{\isachardoublequoteclose}\isanewline
\isanewline
\isacommand{instance}\isamarkupfalse%
\ \isacommand{proof}\isamarkupfalse%
\isanewline
\ \ \isacommand{fix}\isamarkupfalse%
\ i\ j\ k\ {\isacharcolon}{\isacharcolon}\ int\ \isacommand{have}\isamarkupfalse%
\ {\isachardoublequoteopen}{\isacharparenleft}i\ {\isacharplus}\ j{\isacharparenright}\ {\isacharplus}\ k\ {\isacharequal}\ i\ {\isacharplus}\ {\isacharparenleft}j\ {\isacharplus}\ k{\isacharparenright}{\isachardoublequoteclose}\ \isacommand{by}\isamarkupfalse%
\ simp\isanewline
\ \ \isacommand{then}\isamarkupfalse%
\ \isacommand{show}\isamarkupfalse%
\ {\isachardoublequoteopen}{\isacharparenleft}i\ {\isasymotimes}\ j{\isacharparenright}\ {\isasymotimes}\ k\ {\isacharequal}\ i\ {\isasymotimes}\ {\isacharparenleft}j\ {\isasymotimes}\ k{\isacharparenright}{\isachardoublequoteclose}\isanewline
\ \ \ \ \isacommand{unfolding}\isamarkupfalse%
\ mult{\isacharunderscore}int{\isacharunderscore}def\ \isacommand{{\isachardot}}\isamarkupfalse%
\isanewline
\isacommand{qed}\isamarkupfalse%
\isanewline
\isanewline
\isacommand{end}\isamarkupfalse%
%
\endisatagquote
{\isafoldquote}%
%
\isadelimquote
%
\endisadelimquote
%
\begin{isamarkuptext}%
\noindent \hyperlink{command.instantiation}{\mbox{\isa{\isacommand{instantiation}}}} defines class parameters at a
  particular instance using common specification tools (here,
  \hyperlink{command.definition}{\mbox{\isa{\isacommand{definition}}}}).  The concluding \hyperlink{command.instance}{\mbox{\isa{\isacommand{instance}}}} opens a
  proof that the given parameters actually conform to the class
  specification.  Note that the first proof step is the \hyperlink{method.default}{\mbox{\isa{default}}} method, which for such instance proofs maps to the \hyperlink{method.intro-classes}{\mbox{\isa{intro{\isacharunderscore}classes}}} method.  This reduces an instance judgement to the
  relevant primitive proof goals; typically it is the first method
  applied in an instantiation proof.

  From now on, the type-checker will consider \isa{int} as a \isa{semigroup} automatically, i.e.\ any general results are immediately
  available on concrete instances.

  \medskip Another instance of \isa{semigroup} yields the natural
  numbers:%
\end{isamarkuptext}%
\isamarkuptrue%
%
\isadelimquote
%
\endisadelimquote
%
\isatagquote
\isacommand{instantiation}\isamarkupfalse%
\ nat\ {\isacharcolon}{\isacharcolon}\ semigroup\isanewline
\isakeyword{begin}\isanewline
\isanewline
\isacommand{primrec}\isamarkupfalse%
\ mult{\isacharunderscore}nat\ \isakeyword{where}\isanewline
\ \ {\isachardoublequoteopen}{\isacharparenleft}{\isadigit{0}}{\isasymColon}nat{\isacharparenright}\ {\isasymotimes}\ n\ {\isacharequal}\ n{\isachardoublequoteclose}\isanewline
\ \ {\isacharbar}\ {\isachardoublequoteopen}Suc\ m\ {\isasymotimes}\ n\ {\isacharequal}\ Suc\ {\isacharparenleft}m\ {\isasymotimes}\ n{\isacharparenright}{\isachardoublequoteclose}\isanewline
\isanewline
\isacommand{instance}\isamarkupfalse%
\ \isacommand{proof}\isamarkupfalse%
\isanewline
\ \ \isacommand{fix}\isamarkupfalse%
\ m\ n\ q\ {\isacharcolon}{\isacharcolon}\ nat\ \isanewline
\ \ \isacommand{show}\isamarkupfalse%
\ {\isachardoublequoteopen}m\ {\isasymotimes}\ n\ {\isasymotimes}\ q\ {\isacharequal}\ m\ {\isasymotimes}\ {\isacharparenleft}n\ {\isasymotimes}\ q{\isacharparenright}{\isachardoublequoteclose}\isanewline
\ \ \ \ \isacommand{by}\isamarkupfalse%
\ {\isacharparenleft}induct\ m{\isacharparenright}\ auto\isanewline
\isacommand{qed}\isamarkupfalse%
\isanewline
\isanewline
\isacommand{end}\isamarkupfalse%
%
\endisatagquote
{\isafoldquote}%
%
\isadelimquote
%
\endisadelimquote
%
\begin{isamarkuptext}%
\noindent Note the occurence of the name \isa{mult{\isacharunderscore}nat} in the
  primrec declaration; by default, the local name of a class operation
  \isa{f} to be instantiated on type constructor \isa{{\isasymkappa}} is
  mangled as \isa{f{\isacharunderscore}{\isasymkappa}}.  In case of uncertainty, these names may be
  inspected using the \hyperlink{command.print-context}{\mbox{\isa{\isacommand{print{\isacharunderscore}context}}}} command or the
  corresponding ProofGeneral button.%
\end{isamarkuptext}%
\isamarkuptrue%
%
\isamarkupsubsection{Lifting and parametric types%
}
\isamarkuptrue%
%
\begin{isamarkuptext}%
Overloaded definitions given at a class instantiation may include
  recursion over the syntactic structure of types.  As a canonical
  example, we model product semigroups using our simple algebra:%
\end{isamarkuptext}%
\isamarkuptrue%
%
\isadelimquote
%
\endisadelimquote
%
\isatagquote
\isacommand{instantiation}\isamarkupfalse%
\ prod\ {\isacharcolon}{\isacharcolon}\ {\isacharparenleft}semigroup{\isacharcomma}\ semigroup{\isacharparenright}\ semigroup\isanewline
\isakeyword{begin}\isanewline
\isanewline
\isacommand{definition}\isamarkupfalse%
\isanewline
\ \ mult{\isacharunderscore}prod{\isacharunderscore}def{\isacharcolon}\ {\isachardoublequoteopen}p\isactrlisub {\isadigit{1}}\ {\isasymotimes}\ p\isactrlisub {\isadigit{2}}\ {\isacharequal}\ {\isacharparenleft}fst\ p\isactrlisub {\isadigit{1}}\ {\isasymotimes}\ fst\ p\isactrlisub {\isadigit{2}}{\isacharcomma}\ snd\ p\isactrlisub {\isadigit{1}}\ {\isasymotimes}\ snd\ p\isactrlisub {\isadigit{2}}{\isacharparenright}{\isachardoublequoteclose}\isanewline
\isanewline
\isacommand{instance}\isamarkupfalse%
\ \isacommand{proof}\isamarkupfalse%
\isanewline
\ \ \isacommand{fix}\isamarkupfalse%
\ p\isactrlisub {\isadigit{1}}\ p\isactrlisub {\isadigit{2}}\ p\isactrlisub {\isadigit{3}}\ {\isacharcolon}{\isacharcolon}\ {\isachardoublequoteopen}{\isasymalpha}{\isasymColon}semigroup\ {\isasymtimes}\ {\isasymbeta}{\isasymColon}semigroup{\isachardoublequoteclose}\isanewline
\ \ \isacommand{show}\isamarkupfalse%
\ {\isachardoublequoteopen}p\isactrlisub {\isadigit{1}}\ {\isasymotimes}\ p\isactrlisub {\isadigit{2}}\ {\isasymotimes}\ p\isactrlisub {\isadigit{3}}\ {\isacharequal}\ p\isactrlisub {\isadigit{1}}\ {\isasymotimes}\ {\isacharparenleft}p\isactrlisub {\isadigit{2}}\ {\isasymotimes}\ p\isactrlisub {\isadigit{3}}{\isacharparenright}{\isachardoublequoteclose}\isanewline
\ \ \ \ \isacommand{unfolding}\isamarkupfalse%
\ mult{\isacharunderscore}prod{\isacharunderscore}def\ \isacommand{by}\isamarkupfalse%
\ {\isacharparenleft}simp\ add{\isacharcolon}\ assoc{\isacharparenright}\isanewline
\isacommand{qed}\isamarkupfalse%
\ \ \ \ \ \ \isanewline
\isanewline
\isacommand{end}\isamarkupfalse%
%
\endisatagquote
{\isafoldquote}%
%
\isadelimquote
%
\endisadelimquote
%
\begin{isamarkuptext}%
\noindent Associativity of product semigroups is established using
  the definition of \isa{{\isacharparenleft}{\isasymotimes}{\isacharparenright}} on products and the hypothetical
  associativity of the type components; these hypotheses are
  legitimate due to the \isa{semigroup} constraints imposed on the
  type components by the \hyperlink{command.instance}{\mbox{\isa{\isacommand{instance}}}} proposition.  Indeed,
  this pattern often occurs with parametric types and type classes.%
\end{isamarkuptext}%
\isamarkuptrue%
%
\isamarkupsubsection{Subclassing%
}
\isamarkuptrue%
%
\begin{isamarkuptext}%
We define a subclass \isa{monoidl} (a semigroup with a left-hand
  neutral) by extending \isa{semigroup} with one additional
  parameter \isa{neutral} together with its characteristic property:%
\end{isamarkuptext}%
\isamarkuptrue%
%
\isadelimquote
%
\endisadelimquote
%
\isatagquote
\isacommand{class}\isamarkupfalse%
\ monoidl\ {\isacharequal}\ semigroup\ {\isacharplus}\isanewline
\ \ \isakeyword{fixes}\ neutral\ {\isacharcolon}{\isacharcolon}\ {\isachardoublequoteopen}{\isasymalpha}{\isachardoublequoteclose}\ {\isacharparenleft}{\isachardoublequoteopen}{\isasymone}{\isachardoublequoteclose}{\isacharparenright}\isanewline
\ \ \isakeyword{assumes}\ neutl{\isacharcolon}\ {\isachardoublequoteopen}{\isasymone}\ {\isasymotimes}\ x\ {\isacharequal}\ x{\isachardoublequoteclose}%
\endisatagquote
{\isafoldquote}%
%
\isadelimquote
%
\endisadelimquote
%
\begin{isamarkuptext}%
\noindent Again, we prove some instances, by providing suitable
  parameter definitions and proofs for the additional specifications.
  Observe that instantiations for types with the same arity may be
  simultaneous:%
\end{isamarkuptext}%
\isamarkuptrue%
%
\isadelimquote
%
\endisadelimquote
%
\isatagquote
\isacommand{instantiation}\isamarkupfalse%
\ nat\ \isakeyword{and}\ int\ {\isacharcolon}{\isacharcolon}\ monoidl\isanewline
\isakeyword{begin}\isanewline
\isanewline
\isacommand{definition}\isamarkupfalse%
\isanewline
\ \ neutral{\isacharunderscore}nat{\isacharunderscore}def{\isacharcolon}\ {\isachardoublequoteopen}{\isasymone}\ {\isacharequal}\ {\isacharparenleft}{\isadigit{0}}{\isasymColon}nat{\isacharparenright}{\isachardoublequoteclose}\isanewline
\isanewline
\isacommand{definition}\isamarkupfalse%
\isanewline
\ \ neutral{\isacharunderscore}int{\isacharunderscore}def{\isacharcolon}\ {\isachardoublequoteopen}{\isasymone}\ {\isacharequal}\ {\isacharparenleft}{\isadigit{0}}{\isasymColon}int{\isacharparenright}{\isachardoublequoteclose}\isanewline
\isanewline
\isacommand{instance}\isamarkupfalse%
\ \isacommand{proof}\isamarkupfalse%
\isanewline
\ \ \isacommand{fix}\isamarkupfalse%
\ n\ {\isacharcolon}{\isacharcolon}\ nat\isanewline
\ \ \isacommand{show}\isamarkupfalse%
\ {\isachardoublequoteopen}{\isasymone}\ {\isasymotimes}\ n\ {\isacharequal}\ n{\isachardoublequoteclose}\isanewline
\ \ \ \ \isacommand{unfolding}\isamarkupfalse%
\ neutral{\isacharunderscore}nat{\isacharunderscore}def\ \isacommand{by}\isamarkupfalse%
\ simp\isanewline
\isacommand{next}\isamarkupfalse%
\isanewline
\ \ \isacommand{fix}\isamarkupfalse%
\ k\ {\isacharcolon}{\isacharcolon}\ int\isanewline
\ \ \isacommand{show}\isamarkupfalse%
\ {\isachardoublequoteopen}{\isasymone}\ {\isasymotimes}\ k\ {\isacharequal}\ k{\isachardoublequoteclose}\isanewline
\ \ \ \ \isacommand{unfolding}\isamarkupfalse%
\ neutral{\isacharunderscore}int{\isacharunderscore}def\ mult{\isacharunderscore}int{\isacharunderscore}def\ \isacommand{by}\isamarkupfalse%
\ simp\isanewline
\isacommand{qed}\isamarkupfalse%
\isanewline
\isanewline
\isacommand{end}\isamarkupfalse%
\isanewline
\isanewline
\isacommand{instantiation}\isamarkupfalse%
\ prod\ {\isacharcolon}{\isacharcolon}\ {\isacharparenleft}monoidl{\isacharcomma}\ monoidl{\isacharparenright}\ monoidl\isanewline
\isakeyword{begin}\isanewline
\isanewline
\isacommand{definition}\isamarkupfalse%
\isanewline
\ \ neutral{\isacharunderscore}prod{\isacharunderscore}def{\isacharcolon}\ {\isachardoublequoteopen}{\isasymone}\ {\isacharequal}\ {\isacharparenleft}{\isasymone}{\isacharcomma}\ {\isasymone}{\isacharparenright}{\isachardoublequoteclose}\isanewline
\isanewline
\isacommand{instance}\isamarkupfalse%
\ \isacommand{proof}\isamarkupfalse%
\isanewline
\ \ \isacommand{fix}\isamarkupfalse%
\ p\ {\isacharcolon}{\isacharcolon}\ {\isachardoublequoteopen}{\isasymalpha}{\isasymColon}monoidl\ {\isasymtimes}\ {\isasymbeta}{\isasymColon}monoidl{\isachardoublequoteclose}\isanewline
\ \ \isacommand{show}\isamarkupfalse%
\ {\isachardoublequoteopen}{\isasymone}\ {\isasymotimes}\ p\ {\isacharequal}\ p{\isachardoublequoteclose}\isanewline
\ \ \ \ \isacommand{unfolding}\isamarkupfalse%
\ neutral{\isacharunderscore}prod{\isacharunderscore}def\ mult{\isacharunderscore}prod{\isacharunderscore}def\ \isacommand{by}\isamarkupfalse%
\ {\isacharparenleft}simp\ add{\isacharcolon}\ neutl{\isacharparenright}\isanewline
\isacommand{qed}\isamarkupfalse%
\isanewline
\isanewline
\isacommand{end}\isamarkupfalse%
%
\endisatagquote
{\isafoldquote}%
%
\isadelimquote
%
\endisadelimquote
%
\begin{isamarkuptext}%
\noindent Fully-fledged monoids are modelled by another subclass,
  which does not add new parameters but tightens the specification:%
\end{isamarkuptext}%
\isamarkuptrue%
%
\isadelimquote
%
\endisadelimquote
%
\isatagquote
\isacommand{class}\isamarkupfalse%
\ monoid\ {\isacharequal}\ monoidl\ {\isacharplus}\isanewline
\ \ \isakeyword{assumes}\ neutr{\isacharcolon}\ {\isachardoublequoteopen}x\ {\isasymotimes}\ {\isasymone}\ {\isacharequal}\ x{\isachardoublequoteclose}\isanewline
\isanewline
\isacommand{instantiation}\isamarkupfalse%
\ nat\ \isakeyword{and}\ int\ {\isacharcolon}{\isacharcolon}\ monoid\ \isanewline
\isakeyword{begin}\isanewline
\isanewline
\isacommand{instance}\isamarkupfalse%
\ \isacommand{proof}\isamarkupfalse%
\isanewline
\ \ \isacommand{fix}\isamarkupfalse%
\ n\ {\isacharcolon}{\isacharcolon}\ nat\isanewline
\ \ \isacommand{show}\isamarkupfalse%
\ {\isachardoublequoteopen}n\ {\isasymotimes}\ {\isasymone}\ {\isacharequal}\ n{\isachardoublequoteclose}\isanewline
\ \ \ \ \isacommand{unfolding}\isamarkupfalse%
\ neutral{\isacharunderscore}nat{\isacharunderscore}def\ \isacommand{by}\isamarkupfalse%
\ {\isacharparenleft}induct\ n{\isacharparenright}\ simp{\isacharunderscore}all\isanewline
\isacommand{next}\isamarkupfalse%
\isanewline
\ \ \isacommand{fix}\isamarkupfalse%
\ k\ {\isacharcolon}{\isacharcolon}\ int\isanewline
\ \ \isacommand{show}\isamarkupfalse%
\ {\isachardoublequoteopen}k\ {\isasymotimes}\ {\isasymone}\ {\isacharequal}\ k{\isachardoublequoteclose}\isanewline
\ \ \ \ \isacommand{unfolding}\isamarkupfalse%
\ neutral{\isacharunderscore}int{\isacharunderscore}def\ mult{\isacharunderscore}int{\isacharunderscore}def\ \isacommand{by}\isamarkupfalse%
\ simp\isanewline
\isacommand{qed}\isamarkupfalse%
\isanewline
\isanewline
\isacommand{end}\isamarkupfalse%
\isanewline
\isanewline
\isacommand{instantiation}\isamarkupfalse%
\ prod\ {\isacharcolon}{\isacharcolon}\ {\isacharparenleft}monoid{\isacharcomma}\ monoid{\isacharparenright}\ monoid\isanewline
\isakeyword{begin}\isanewline
\isanewline
\isacommand{instance}\isamarkupfalse%
\ \isacommand{proof}\isamarkupfalse%
\ \isanewline
\ \ \isacommand{fix}\isamarkupfalse%
\ p\ {\isacharcolon}{\isacharcolon}\ {\isachardoublequoteopen}{\isasymalpha}{\isasymColon}monoid\ {\isasymtimes}\ {\isasymbeta}{\isasymColon}monoid{\isachardoublequoteclose}\isanewline
\ \ \isacommand{show}\isamarkupfalse%
\ {\isachardoublequoteopen}p\ {\isasymotimes}\ {\isasymone}\ {\isacharequal}\ p{\isachardoublequoteclose}\isanewline
\ \ \ \ \isacommand{unfolding}\isamarkupfalse%
\ neutral{\isacharunderscore}prod{\isacharunderscore}def\ mult{\isacharunderscore}prod{\isacharunderscore}def\ \isacommand{by}\isamarkupfalse%
\ {\isacharparenleft}simp\ add{\isacharcolon}\ neutr{\isacharparenright}\isanewline
\isacommand{qed}\isamarkupfalse%
\isanewline
\isanewline
\isacommand{end}\isamarkupfalse%
%
\endisatagquote
{\isafoldquote}%
%
\isadelimquote
%
\endisadelimquote
%
\begin{isamarkuptext}%
\noindent To finish our small algebra example, we add a \isa{group} class with a corresponding instance:%
\end{isamarkuptext}%
\isamarkuptrue%
%
\isadelimquote
%
\endisadelimquote
%
\isatagquote
\isacommand{class}\isamarkupfalse%
\ group\ {\isacharequal}\ monoidl\ {\isacharplus}\isanewline
\ \ \isakeyword{fixes}\ inverse\ {\isacharcolon}{\isacharcolon}\ {\isachardoublequoteopen}{\isasymalpha}\ {\isasymRightarrow}\ {\isasymalpha}{\isachardoublequoteclose}\ \ \ \ {\isacharparenleft}{\isachardoublequoteopen}{\isacharparenleft}{\isacharunderscore}{\isasymdiv}{\isacharparenright}{\isachardoublequoteclose}\ {\isacharbrackleft}{\isadigit{1}}{\isadigit{0}}{\isadigit{0}}{\isadigit{0}}{\isacharbrackright}\ {\isadigit{9}}{\isadigit{9}}{\isadigit{9}}{\isacharparenright}\isanewline
\ \ \isakeyword{assumes}\ invl{\isacharcolon}\ {\isachardoublequoteopen}x{\isasymdiv}\ {\isasymotimes}\ x\ {\isacharequal}\ {\isasymone}{\isachardoublequoteclose}\isanewline
\isanewline
\isacommand{instantiation}\isamarkupfalse%
\ int\ {\isacharcolon}{\isacharcolon}\ group\isanewline
\isakeyword{begin}\isanewline
\isanewline
\isacommand{definition}\isamarkupfalse%
\isanewline
\ \ inverse{\isacharunderscore}int{\isacharunderscore}def{\isacharcolon}\ {\isachardoublequoteopen}i{\isasymdiv}\ {\isacharequal}\ {\isacharminus}\ {\isacharparenleft}i{\isasymColon}int{\isacharparenright}{\isachardoublequoteclose}\isanewline
\isanewline
\isacommand{instance}\isamarkupfalse%
\ \isacommand{proof}\isamarkupfalse%
\isanewline
\ \ \isacommand{fix}\isamarkupfalse%
\ i\ {\isacharcolon}{\isacharcolon}\ int\isanewline
\ \ \isacommand{have}\isamarkupfalse%
\ {\isachardoublequoteopen}{\isacharminus}i\ {\isacharplus}\ i\ {\isacharequal}\ {\isadigit{0}}{\isachardoublequoteclose}\ \isacommand{by}\isamarkupfalse%
\ simp\isanewline
\ \ \isacommand{then}\isamarkupfalse%
\ \isacommand{show}\isamarkupfalse%
\ {\isachardoublequoteopen}i{\isasymdiv}\ {\isasymotimes}\ i\ {\isacharequal}\ {\isasymone}{\isachardoublequoteclose}\isanewline
\ \ \ \ \isacommand{unfolding}\isamarkupfalse%
\ mult{\isacharunderscore}int{\isacharunderscore}def\ neutral{\isacharunderscore}int{\isacharunderscore}def\ inverse{\isacharunderscore}int{\isacharunderscore}def\ \isacommand{{\isachardot}}\isamarkupfalse%
\isanewline
\isacommand{qed}\isamarkupfalse%
\isanewline
\isanewline
\isacommand{end}\isamarkupfalse%
%
\endisatagquote
{\isafoldquote}%
%
\isadelimquote
%
\endisadelimquote
%
\isamarkupsection{Type classes as locales%
}
\isamarkuptrue%
%
\isamarkupsubsection{A look behind the scenes%
}
\isamarkuptrue%
%
\begin{isamarkuptext}%
The example above gives an impression how Isar type classes work in
  practice.  As stated in the introduction, classes also provide a
  link to Isar's locale system.  Indeed, the logical core of a class
  is nothing other than a locale:%
\end{isamarkuptext}%
\isamarkuptrue%
%
\isadelimquote
%
\endisadelimquote
%
\isatagquote
\isacommand{class}\isamarkupfalse%
\ idem\ {\isacharequal}\isanewline
\ \ \isakeyword{fixes}\ f\ {\isacharcolon}{\isacharcolon}\ {\isachardoublequoteopen}{\isasymalpha}\ {\isasymRightarrow}\ {\isasymalpha}{\isachardoublequoteclose}\isanewline
\ \ \isakeyword{assumes}\ idem{\isacharcolon}\ {\isachardoublequoteopen}f\ {\isacharparenleft}f\ x{\isacharparenright}\ {\isacharequal}\ f\ x{\isachardoublequoteclose}%
\endisatagquote
{\isafoldquote}%
%
\isadelimquote
%
\endisadelimquote
%
\begin{isamarkuptext}%
\noindent essentially introduces the locale%
\end{isamarkuptext}%
\isamarkuptrue%
\ %
\isadeliminvisible
%
\endisadeliminvisible
%
\isataginvisible
%
\endisataginvisible
{\isafoldinvisible}%
%
\isadeliminvisible
%
\endisadeliminvisible
%
\isadelimquote
%
\endisadelimquote
%
\isatagquote
\isacommand{locale}\isamarkupfalse%
\ idem\ {\isacharequal}\isanewline
\ \ \isakeyword{fixes}\ f\ {\isacharcolon}{\isacharcolon}\ {\isachardoublequoteopen}{\isasymalpha}\ {\isasymRightarrow}\ {\isasymalpha}{\isachardoublequoteclose}\isanewline
\ \ \isakeyword{assumes}\ idem{\isacharcolon}\ {\isachardoublequoteopen}f\ {\isacharparenleft}f\ x{\isacharparenright}\ {\isacharequal}\ f\ x{\isachardoublequoteclose}%
\endisatagquote
{\isafoldquote}%
%
\isadelimquote
%
\endisadelimquote
%
\begin{isamarkuptext}%
\noindent together with corresponding constant(s):%
\end{isamarkuptext}%
\isamarkuptrue%
%
\isadelimquote
%
\endisadelimquote
%
\isatagquote
\isacommand{consts}\isamarkupfalse%
\ f\ {\isacharcolon}{\isacharcolon}\ {\isachardoublequoteopen}{\isasymalpha}\ {\isasymRightarrow}\ {\isasymalpha}{\isachardoublequoteclose}%
\endisatagquote
{\isafoldquote}%
%
\isadelimquote
%
\endisadelimquote
%
\begin{isamarkuptext}%
\noindent The connection to the type system is done by means
  of a primitive type class%
\end{isamarkuptext}%
\isamarkuptrue%
\ %
\isadeliminvisible
%
\endisadeliminvisible
%
\isataginvisible
%
\endisataginvisible
{\isafoldinvisible}%
%
\isadeliminvisible
%
\endisadeliminvisible
%
\isadelimquote
%
\endisadelimquote
%
\isatagquote
\isacommand{classes}\isamarkupfalse%
\ idem\ {\isacharless}\ type%
\endisatagquote
{\isafoldquote}%
%
\isadelimquote
%
\endisadelimquote
%
\isadeliminvisible
%
\endisadeliminvisible
%
\isataginvisible
%
\endisataginvisible
{\isafoldinvisible}%
%
\isadeliminvisible
%
\endisadeliminvisible
%
\begin{isamarkuptext}%
\noindent together with a corresponding interpretation:%
\end{isamarkuptext}%
\isamarkuptrue%
%
\isadelimquote
%
\endisadelimquote
%
\isatagquote
\isacommand{interpretation}\isamarkupfalse%
\ idem{\isacharunderscore}class{\isacharcolon}\isanewline
\ \ idem\ {\isachardoublequoteopen}f\ {\isasymColon}\ {\isacharparenleft}{\isasymalpha}{\isasymColon}idem{\isacharparenright}\ {\isasymRightarrow}\ {\isasymalpha}{\isachardoublequoteclose}%
\endisatagquote
{\isafoldquote}%
%
\isadelimquote
%
\endisadelimquote
%
\begin{isamarkuptext}%
\noindent This gives you the full power of the Isabelle module system;
  conclusions in locale \isa{idem} are implicitly propagated
  to class \isa{idem}.%
\end{isamarkuptext}%
\isamarkuptrue%
\ %
\isadeliminvisible
%
\endisadeliminvisible
%
\isataginvisible
%
\endisataginvisible
{\isafoldinvisible}%
%
\isadeliminvisible
%
\endisadeliminvisible
%
\isamarkupsubsection{Abstract reasoning%
}
\isamarkuptrue%
%
\begin{isamarkuptext}%
Isabelle locales enable reasoning at a general level, while results
  are implicitly transferred to all instances.  For example, we can
  now establish the \isa{left{\isacharunderscore}cancel} lemma for groups, which
  states that the function \isa{{\isacharparenleft}x\ {\isasymotimes}{\isacharparenright}} is injective:%
\end{isamarkuptext}%
\isamarkuptrue%
%
\isadelimquote
%
\endisadelimquote
%
\isatagquote
\isacommand{lemma}\isamarkupfalse%
\ {\isacharparenleft}\isakeyword{in}\ group{\isacharparenright}\ left{\isacharunderscore}cancel{\isacharcolon}\ {\isachardoublequoteopen}x\ {\isasymotimes}\ y\ {\isacharequal}\ x\ {\isasymotimes}\ z\ {\isasymlongleftrightarrow}\ y\ {\isacharequal}\ z{\isachardoublequoteclose}\isanewline
\isacommand{proof}\isamarkupfalse%
\isanewline
\ \ \isacommand{assume}\isamarkupfalse%
\ {\isachardoublequoteopen}x\ {\isasymotimes}\ y\ {\isacharequal}\ x\ {\isasymotimes}\ z{\isachardoublequoteclose}\isanewline
\ \ \isacommand{then}\isamarkupfalse%
\ \isacommand{have}\isamarkupfalse%
\ {\isachardoublequoteopen}x{\isasymdiv}\ {\isasymotimes}\ {\isacharparenleft}x\ {\isasymotimes}\ y{\isacharparenright}\ {\isacharequal}\ x{\isasymdiv}\ {\isasymotimes}\ {\isacharparenleft}x\ {\isasymotimes}\ z{\isacharparenright}{\isachardoublequoteclose}\ \isacommand{by}\isamarkupfalse%
\ simp\isanewline
\ \ \isacommand{then}\isamarkupfalse%
\ \isacommand{have}\isamarkupfalse%
\ {\isachardoublequoteopen}{\isacharparenleft}x{\isasymdiv}\ {\isasymotimes}\ x{\isacharparenright}\ {\isasymotimes}\ y\ {\isacharequal}\ {\isacharparenleft}x{\isasymdiv}\ {\isasymotimes}\ x{\isacharparenright}\ {\isasymotimes}\ z{\isachardoublequoteclose}\ \isacommand{using}\isamarkupfalse%
\ assoc\ \isacommand{by}\isamarkupfalse%
\ simp\isanewline
\ \ \isacommand{then}\isamarkupfalse%
\ \isacommand{show}\isamarkupfalse%
\ {\isachardoublequoteopen}y\ {\isacharequal}\ z{\isachardoublequoteclose}\ \isacommand{using}\isamarkupfalse%
\ neutl\ \isakeyword{and}\ invl\ \isacommand{by}\isamarkupfalse%
\ simp\isanewline
\isacommand{next}\isamarkupfalse%
\isanewline
\ \ \isacommand{assume}\isamarkupfalse%
\ {\isachardoublequoteopen}y\ {\isacharequal}\ z{\isachardoublequoteclose}\isanewline
\ \ \isacommand{then}\isamarkupfalse%
\ \isacommand{show}\isamarkupfalse%
\ {\isachardoublequoteopen}x\ {\isasymotimes}\ y\ {\isacharequal}\ x\ {\isasymotimes}\ z{\isachardoublequoteclose}\ \isacommand{by}\isamarkupfalse%
\ simp\isanewline
\isacommand{qed}\isamarkupfalse%
%
\endisatagquote
{\isafoldquote}%
%
\isadelimquote
%
\endisadelimquote
%
\begin{isamarkuptext}%
\noindent Here the \qt{\hyperlink{keyword.in}{\mbox{\isa{\isakeyword{in}}}} \isa{group}} target
  specification indicates that the result is recorded within that
  context for later use.  This local theorem is also lifted to the
  global one \hyperlink{fact.group.left-cancel:}{\mbox{\isa{group{\isachardot}left{\isacharunderscore}cancel{\isacharcolon}}}} \isa{{\isachardoublequote}{\isasymAnd}x\ y\ z\ {\isasymColon}\ {\isasymalpha}{\isasymColon}group{\isachardot}\ x\ {\isasymotimes}\ y\ {\isacharequal}\ x\ {\isasymotimes}\ z\ {\isasymlongleftrightarrow}\ y\ {\isacharequal}\ z{\isachardoublequote}}.  Since type \isa{int} has been
  made an instance of \isa{group} before, we may refer to that
  fact as well: \isa{{\isachardoublequote}{\isasymAnd}x\ y\ z\ {\isasymColon}\ int{\isachardot}\ x\ {\isasymotimes}\ y\ {\isacharequal}\ x\ {\isasymotimes}\ z\ {\isasymlongleftrightarrow}\ y\ {\isacharequal}\ z{\isachardoublequote}}.%
\end{isamarkuptext}%
\isamarkuptrue%
%
\isamarkupsubsection{Derived definitions%
}
\isamarkuptrue%
%
\begin{isamarkuptext}%
Isabelle locales are targets which support local definitions:%
\end{isamarkuptext}%
\isamarkuptrue%
%
\isadelimquote
%
\endisadelimquote
%
\isatagquote
\isacommand{primrec}\isamarkupfalse%
\ {\isacharparenleft}\isakeyword{in}\ monoid{\isacharparenright}\ pow{\isacharunderscore}nat\ {\isacharcolon}{\isacharcolon}\ {\isachardoublequoteopen}nat\ {\isasymRightarrow}\ {\isasymalpha}\ {\isasymRightarrow}\ {\isasymalpha}{\isachardoublequoteclose}\ \isakeyword{where}\isanewline
\ \ {\isachardoublequoteopen}pow{\isacharunderscore}nat\ {\isadigit{0}}\ x\ {\isacharequal}\ {\isasymone}{\isachardoublequoteclose}\isanewline
\ \ {\isacharbar}\ {\isachardoublequoteopen}pow{\isacharunderscore}nat\ {\isacharparenleft}Suc\ n{\isacharparenright}\ x\ {\isacharequal}\ x\ {\isasymotimes}\ pow{\isacharunderscore}nat\ n\ x{\isachardoublequoteclose}%
\endisatagquote
{\isafoldquote}%
%
\isadelimquote
%
\endisadelimquote
%
\begin{isamarkuptext}%
\noindent If the locale \isa{group} is also a class, this local
  definition is propagated onto a global definition of \isa{{\isachardoublequote}pow{\isacharunderscore}nat\ {\isasymColon}\ nat\ {\isasymRightarrow}\ {\isasymalpha}{\isasymColon}monoid\ {\isasymRightarrow}\ {\isasymalpha}{\isasymColon}monoid{\isachardoublequote}} with corresponding theorems

  \isa{pow{\isacharunderscore}nat\ {\isadigit{0}}\ x\ {\isacharequal}\ {\isasymone}\isasep\isanewline%
pow{\isacharunderscore}nat\ {\isacharparenleft}Suc\ n{\isacharparenright}\ x\ {\isacharequal}\ x\ {\isasymotimes}\ pow{\isacharunderscore}nat\ n\ x}.

  \noindent As you can see from this example, for local definitions
  you may use any specification tool which works together with
  locales, such as Krauss's recursive function package
  \cite{krauss2006}.%
\end{isamarkuptext}%
\isamarkuptrue%
%
\isamarkupsubsection{A functor analogy%
}
\isamarkuptrue%
%
\begin{isamarkuptext}%
We introduced Isar classes by analogy to type classes in functional
  programming; if we reconsider this in the context of what has been
  said about type classes and locales, we can drive this analogy
  further by stating that type classes essentially correspond to
  functors that have a canonical interpretation as type classes.
  There is also the possibility of other interpretations.  For
  example, \isa{list}s also form a monoid with \isa{append} and
  \isa{{\isacharbrackleft}{\isacharbrackright}} as operations, but it seems inappropriate to apply to
  lists the same operations as for genuinely algebraic types.  In such
  a case, we can simply make a particular interpretation of monoids
  for lists:%
\end{isamarkuptext}%
\isamarkuptrue%
%
\isadelimquote
%
\endisadelimquote
%
\isatagquote
\isacommand{interpretation}\isamarkupfalse%
\ list{\isacharunderscore}monoid{\isacharcolon}\ monoid\ append\ {\isachardoublequoteopen}{\isacharbrackleft}{\isacharbrackright}{\isachardoublequoteclose}\isanewline
\ \ \isacommand{proof}\isamarkupfalse%
\ \isacommand{qed}\isamarkupfalse%
\ auto%
\endisatagquote
{\isafoldquote}%
%
\isadelimquote
%
\endisadelimquote
%
\begin{isamarkuptext}%
\noindent This enables us to apply facts on monoids
  to lists, e.g. \isa{{\isacharbrackleft}{\isacharbrackright}\ {\isacharat}\ x\ {\isacharequal}\ x}.

  When using this interpretation pattern, it may also
  be appropriate to map derived definitions accordingly:%
\end{isamarkuptext}%
\isamarkuptrue%
%
\isadelimquote
%
\endisadelimquote
%
\isatagquote
\isacommand{primrec}\isamarkupfalse%
\ replicate\ {\isacharcolon}{\isacharcolon}\ {\isachardoublequoteopen}nat\ {\isasymRightarrow}\ {\isasymalpha}\ list\ {\isasymRightarrow}\ {\isasymalpha}\ list{\isachardoublequoteclose}\ \isakeyword{where}\isanewline
\ \ {\isachardoublequoteopen}replicate\ {\isadigit{0}}\ {\isacharunderscore}\ {\isacharequal}\ {\isacharbrackleft}{\isacharbrackright}{\isachardoublequoteclose}\isanewline
\ \ {\isacharbar}\ {\isachardoublequoteopen}replicate\ {\isacharparenleft}Suc\ n{\isacharparenright}\ xs\ {\isacharequal}\ xs\ {\isacharat}\ replicate\ n\ xs{\isachardoublequoteclose}\isanewline
\isanewline
\isacommand{interpretation}\isamarkupfalse%
\ list{\isacharunderscore}monoid{\isacharcolon}\ monoid\ append\ {\isachardoublequoteopen}{\isacharbrackleft}{\isacharbrackright}{\isachardoublequoteclose}\ \isakeyword{where}\isanewline
\ \ {\isachardoublequoteopen}monoid{\isachardot}pow{\isacharunderscore}nat\ append\ {\isacharbrackleft}{\isacharbrackright}\ {\isacharequal}\ replicate{\isachardoublequoteclose}\isanewline
\isacommand{proof}\isamarkupfalse%
\ {\isacharminus}\isanewline
\ \ \isacommand{interpret}\isamarkupfalse%
\ monoid\ append\ {\isachardoublequoteopen}{\isacharbrackleft}{\isacharbrackright}{\isachardoublequoteclose}\ \isacommand{{\isachardot}{\isachardot}}\isamarkupfalse%
\isanewline
\ \ \isacommand{show}\isamarkupfalse%
\ {\isachardoublequoteopen}monoid{\isachardot}pow{\isacharunderscore}nat\ append\ {\isacharbrackleft}{\isacharbrackright}\ {\isacharequal}\ replicate{\isachardoublequoteclose}\isanewline
\ \ \isacommand{proof}\isamarkupfalse%
\isanewline
\ \ \ \ \isacommand{fix}\isamarkupfalse%
\ n\isanewline
\ \ \ \ \isacommand{show}\isamarkupfalse%
\ {\isachardoublequoteopen}monoid{\isachardot}pow{\isacharunderscore}nat\ append\ {\isacharbrackleft}{\isacharbrackright}\ n\ {\isacharequal}\ replicate\ n{\isachardoublequoteclose}\isanewline
\ \ \ \ \ \ \isacommand{by}\isamarkupfalse%
\ {\isacharparenleft}induct\ n{\isacharparenright}\ auto\isanewline
\ \ \isacommand{qed}\isamarkupfalse%
\isanewline
\isacommand{qed}\isamarkupfalse%
\ intro{\isacharunderscore}locales%
\endisatagquote
{\isafoldquote}%
%
\isadelimquote
%
\endisadelimquote
%
\begin{isamarkuptext}%
\noindent This pattern is also helpful to reuse abstract
  specifications on the \emph{same} type.  For example, think of a
  class \isa{preorder}; for type \isa{nat}, there are at least two
  possible instances: the natural order or the order induced by the
  divides relation.  But only one of these instances can be used for
  \hyperlink{command.instantiation}{\mbox{\isa{\isacommand{instantiation}}}}; using the locale behind the class \isa{preorder}, it is still possible to utilise the same abstract
  specification again using \hyperlink{command.interpretation}{\mbox{\isa{\isacommand{interpretation}}}}.%
\end{isamarkuptext}%
\isamarkuptrue%
%
\isamarkupsubsection{Additional subclass relations%
}
\isamarkuptrue%
%
\begin{isamarkuptext}%
Any \isa{group} is also a \isa{monoid}; this can be made
  explicit by claiming an additional subclass relation, together with
  a proof of the logical difference:%
\end{isamarkuptext}%
\isamarkuptrue%
%
\isadelimquote
%
\endisadelimquote
%
\isatagquote
\isacommand{subclass}\isamarkupfalse%
\ {\isacharparenleft}\isakeyword{in}\ group{\isacharparenright}\ monoid\isanewline
\isacommand{proof}\isamarkupfalse%
\isanewline
\ \ \isacommand{fix}\isamarkupfalse%
\ x\isanewline
\ \ \isacommand{from}\isamarkupfalse%
\ invl\ \isacommand{have}\isamarkupfalse%
\ {\isachardoublequoteopen}x{\isasymdiv}\ {\isasymotimes}\ x\ {\isacharequal}\ {\isasymone}{\isachardoublequoteclose}\ \isacommand{by}\isamarkupfalse%
\ simp\isanewline
\ \ \isacommand{with}\isamarkupfalse%
\ assoc\ {\isacharbrackleft}symmetric{\isacharbrackright}\ neutl\ invl\ \isacommand{have}\isamarkupfalse%
\ {\isachardoublequoteopen}x{\isasymdiv}\ {\isasymotimes}\ {\isacharparenleft}x\ {\isasymotimes}\ {\isasymone}{\isacharparenright}\ {\isacharequal}\ x{\isasymdiv}\ {\isasymotimes}\ x{\isachardoublequoteclose}\ \isacommand{by}\isamarkupfalse%
\ simp\isanewline
\ \ \isacommand{with}\isamarkupfalse%
\ left{\isacharunderscore}cancel\ \isacommand{show}\isamarkupfalse%
\ {\isachardoublequoteopen}x\ {\isasymotimes}\ {\isasymone}\ {\isacharequal}\ x{\isachardoublequoteclose}\ \isacommand{by}\isamarkupfalse%
\ simp\isanewline
\isacommand{qed}\isamarkupfalse%
%
\endisatagquote
{\isafoldquote}%
%
\isadelimquote
%
\endisadelimquote
%
\begin{isamarkuptext}%
The logical proof is carried out on the locale level.  Afterwards it
  is propagated to the type system, making \isa{group} an instance
  of \isa{monoid} by adding an additional edge to the graph of
  subclass relations (\figref{fig:subclass}).

  \begin{figure}[htbp]
   \begin{center}
     \small
     \unitlength 0.6mm
     \begin{picture}(40,60)(0,0)
       \put(20,60){\makebox(0,0){\isa{semigroup}}}
       \put(20,40){\makebox(0,0){\isa{monoidl}}}
       \put(00,20){\makebox(0,0){\isa{monoid}}}
       \put(40,00){\makebox(0,0){\isa{group}}}
       \put(20,55){\vector(0,-1){10}}
       \put(15,35){\vector(-1,-1){10}}
       \put(25,35){\vector(1,-3){10}}
     \end{picture}
     \hspace{8em}
     \begin{picture}(40,60)(0,0)
       \put(20,60){\makebox(0,0){\isa{semigroup}}}
       \put(20,40){\makebox(0,0){\isa{monoidl}}}
       \put(00,20){\makebox(0,0){\isa{monoid}}}
       \put(40,00){\makebox(0,0){\isa{group}}}
       \put(20,55){\vector(0,-1){10}}
       \put(15,35){\vector(-1,-1){10}}
       \put(05,15){\vector(3,-1){30}}
     \end{picture}
     \caption{Subclass relationship of monoids and groups:
        before and after establishing the relationship
        \isa{group\ {\isasymsubseteq}\ monoid};  transitive edges are left out.}
     \label{fig:subclass}
   \end{center}
  \end{figure}

  For illustration, a derived definition in \isa{group} using \isa{pow{\isacharunderscore}nat}%
\end{isamarkuptext}%
\isamarkuptrue%
%
\isadelimquote
%
\endisadelimquote
%
\isatagquote
\isacommand{definition}\isamarkupfalse%
\ {\isacharparenleft}\isakeyword{in}\ group{\isacharparenright}\ pow{\isacharunderscore}int\ {\isacharcolon}{\isacharcolon}\ {\isachardoublequoteopen}int\ {\isasymRightarrow}\ {\isasymalpha}\ {\isasymRightarrow}\ {\isasymalpha}{\isachardoublequoteclose}\ \isakeyword{where}\isanewline
\ \ {\isachardoublequoteopen}pow{\isacharunderscore}int\ k\ x\ {\isacharequal}\ {\isacharparenleft}if\ k\ {\isachargreater}{\isacharequal}\ {\isadigit{0}}\isanewline
\ \ \ \ then\ pow{\isacharunderscore}nat\ {\isacharparenleft}nat\ k{\isacharparenright}\ x\isanewline
\ \ \ \ else\ {\isacharparenleft}pow{\isacharunderscore}nat\ {\isacharparenleft}nat\ {\isacharparenleft}{\isacharminus}\ k{\isacharparenright}{\isacharparenright}\ x{\isacharparenright}{\isasymdiv}{\isacharparenright}{\isachardoublequoteclose}%
\endisatagquote
{\isafoldquote}%
%
\isadelimquote
%
\endisadelimquote
%
\begin{isamarkuptext}%
\noindent yields the global definition of \isa{{\isachardoublequote}pow{\isacharunderscore}int\ {\isasymColon}\ int\ {\isasymRightarrow}\ {\isasymalpha}{\isasymColon}group\ {\isasymRightarrow}\ {\isasymalpha}{\isasymColon}group{\isachardoublequote}} with the corresponding theorem \isa{pow{\isacharunderscore}int\ k\ x\ {\isacharequal}\ {\isacharparenleft}if\ {\isadigit{0}}\ {\isasymle}\ k\ then\ pow{\isacharunderscore}nat\ {\isacharparenleft}nat\ k{\isacharparenright}\ x\ else\ {\isacharparenleft}pow{\isacharunderscore}nat\ {\isacharparenleft}nat\ {\isacharparenleft}{\isacharminus}\ k{\isacharparenright}{\isacharparenright}\ x{\isacharparenright}{\isasymdiv}{\isacharparenright}}.%
\end{isamarkuptext}%
\isamarkuptrue%
%
\isamarkupsubsection{A note on syntax%
}
\isamarkuptrue%
%
\begin{isamarkuptext}%
As a convenience, class context syntax allows references to local
  class operations and their global counterparts uniformly; type
  inference resolves ambiguities.  For example:%
\end{isamarkuptext}%
\isamarkuptrue%
%
\isadelimquote
%
\endisadelimquote
%
\isatagquote
\isacommand{context}\isamarkupfalse%
\ semigroup\isanewline
\isakeyword{begin}\isanewline
\isanewline
\isacommand{term}\isamarkupfalse%
\ {\isachardoublequoteopen}x\ {\isasymotimes}\ y{\isachardoublequoteclose}\ %
\isamarkupcmt{example 1%
}
\isanewline
\isacommand{term}\isamarkupfalse%
\ {\isachardoublequoteopen}{\isacharparenleft}x{\isasymColon}nat{\isacharparenright}\ {\isasymotimes}\ y{\isachardoublequoteclose}\ %
\isamarkupcmt{example 2%
}
\isanewline
\isanewline
\isacommand{end}\isamarkupfalse%
\isanewline
\isanewline
\isacommand{term}\isamarkupfalse%
\ {\isachardoublequoteopen}x\ {\isasymotimes}\ y{\isachardoublequoteclose}\ %
\isamarkupcmt{example 3%
}
%
\endisatagquote
{\isafoldquote}%
%
\isadelimquote
%
\endisadelimquote
%
\begin{isamarkuptext}%
\noindent Here in example 1, the term refers to the local class
  operation \isa{mult\ {\isacharbrackleft}{\isasymalpha}{\isacharbrackright}}, whereas in example 2 the type
  constraint enforces the global class operation \isa{mult\ {\isacharbrackleft}nat{\isacharbrackright}}.
  In the global context in example 3, the reference is to the
  polymorphic global class operation \isa{mult\ {\isacharbrackleft}{\isacharquery}{\isasymalpha}\ {\isasymColon}\ semigroup{\isacharbrackright}}.%
\end{isamarkuptext}%
\isamarkuptrue%
%
\isamarkupsection{Further issues%
}
\isamarkuptrue%
%
\isamarkupsubsection{Type classes and code generation%
}
\isamarkuptrue%
%
\begin{isamarkuptext}%
Turning back to the first motivation for type classes, namely
  overloading, it is obvious that overloading stemming from \hyperlink{command.class}{\mbox{\isa{\isacommand{class}}}} statements and \hyperlink{command.instantiation}{\mbox{\isa{\isacommand{instantiation}}}} targets naturally
  maps to Haskell type classes.  The code generator framework
  \cite{isabelle-codegen} takes this into account.  If the target
  language (e.g.~SML) lacks type classes, then they are implemented by
  an explicit dictionary construction.  As example, let's go back to
  the power function:%
\end{isamarkuptext}%
\isamarkuptrue%
%
\isadelimquote
%
\endisadelimquote
%
\isatagquote
\isacommand{definition}\isamarkupfalse%
\ example\ {\isacharcolon}{\isacharcolon}\ int\ \isakeyword{where}\isanewline
\ \ {\isachardoublequoteopen}example\ {\isacharequal}\ pow{\isacharunderscore}int\ {\isadigit{1}}{\isadigit{0}}\ {\isacharparenleft}{\isacharminus}{\isadigit{2}}{\isacharparenright}{\isachardoublequoteclose}%
\endisatagquote
{\isafoldquote}%
%
\isadelimquote
%
\endisadelimquote
%
\begin{isamarkuptext}%
\noindent This maps to Haskell as follows:%
\end{isamarkuptext}%
\isamarkuptrue%
%
\isadeliminvisible
%
\endisadeliminvisible
%
\isataginvisible
%
\endisataginvisible
{\isafoldinvisible}%
%
\isadeliminvisible
%
\endisadeliminvisible
%
\isadelimquote
%
\endisadelimquote
%
\isatagquote
%
\begin{isamarkuptext}%
\isatypewriter%
\noindent%
\hspace*{0pt}module Example where {\char123}\\
\hspace*{0pt}\\
\hspace*{0pt}data Nat = Zero{\char95}nat | Suc Nat;\\
\hspace*{0pt}\\
\hspace*{0pt}nat{\char95}aux ::~Integer -> Nat -> Nat;\\
\hspace*{0pt}nat{\char95}aux i n = (if i <= 0 then n else nat{\char95}aux (i - 1) (Suc n));\\
\hspace*{0pt}\\
\hspace*{0pt}nat ::~Integer -> Nat;\\
\hspace*{0pt}nat i = nat{\char95}aux i Zero{\char95}nat;\\
\hspace*{0pt}\\
\hspace*{0pt}class Semigroup a where {\char123}\\
\hspace*{0pt} ~mult ::~a -> a -> a;\\
\hspace*{0pt}{\char125};\\
\hspace*{0pt}\\
\hspace*{0pt}class (Semigroup a) => Monoidl a where {\char123}\\
\hspace*{0pt} ~neutral ::~a;\\
\hspace*{0pt}{\char125};\\
\hspace*{0pt}\\
\hspace*{0pt}class (Monoidl a) => Monoid a where {\char123}\\
\hspace*{0pt}{\char125};\\
\hspace*{0pt}\\
\hspace*{0pt}class (Monoid a) => Group a where {\char123}\\
\hspace*{0pt} ~inverse ::~a -> a;\\
\hspace*{0pt}{\char125};\\
\hspace*{0pt}\\
\hspace*{0pt}pow{\char95}nat ::~forall a.~(Monoid a) => Nat -> a -> a;\\
\hspace*{0pt}pow{\char95}nat Zero{\char95}nat x = neutral;\\
\hspace*{0pt}pow{\char95}nat (Suc n) x = mult x (pow{\char95}nat n x);\\
\hspace*{0pt}\\
\hspace*{0pt}pow{\char95}int ::~forall a.~(Group a) => Integer -> a -> a;\\
\hspace*{0pt}pow{\char95}int k x =\\
\hspace*{0pt} ~(if 0 <= k then pow{\char95}nat (nat k) x\\
\hspace*{0pt} ~~~else inverse (pow{\char95}nat (nat (negate k)) x));\\
\hspace*{0pt}\\
\hspace*{0pt}mult{\char95}int ::~Integer -> Integer -> Integer;\\
\hspace*{0pt}mult{\char95}int i j = i + j;\\
\hspace*{0pt}\\
\hspace*{0pt}instance Semigroup Integer where {\char123}\\
\hspace*{0pt} ~mult = mult{\char95}int;\\
\hspace*{0pt}{\char125};\\
\hspace*{0pt}\\
\hspace*{0pt}neutral{\char95}int ::~Integer;\\
\hspace*{0pt}neutral{\char95}int = 0;\\
\hspace*{0pt}\\
\hspace*{0pt}instance Monoidl Integer where {\char123}\\
\hspace*{0pt} ~neutral = neutral{\char95}int;\\
\hspace*{0pt}{\char125};\\
\hspace*{0pt}\\
\hspace*{0pt}instance Monoid Integer where {\char123}\\
\hspace*{0pt}{\char125};\\
\hspace*{0pt}\\
\hspace*{0pt}inverse{\char95}int ::~Integer -> Integer;\\
\hspace*{0pt}inverse{\char95}int i = negate i;\\
\hspace*{0pt}\\
\hspace*{0pt}instance Group Integer where {\char123}\\
\hspace*{0pt} ~inverse = inverse{\char95}int;\\
\hspace*{0pt}{\char125};\\
\hspace*{0pt}\\
\hspace*{0pt}example ::~Integer;\\
\hspace*{0pt}example = pow{\char95}int 10 (-2);\\
\hspace*{0pt}\\
\hspace*{0pt}{\char125}%
\end{isamarkuptext}%
\isamarkuptrue%
%
\endisatagquote
{\isafoldquote}%
%
\isadelimquote
%
\endisadelimquote
%
\begin{isamarkuptext}%
\noindent The code in SML has explicit dictionary passing:%
\end{isamarkuptext}%
\isamarkuptrue%
%
\isadelimquote
%
\endisadelimquote
%
\isatagquote
%
\begin{isamarkuptext}%
\isatypewriter%
\noindent%
\hspace*{0pt}structure Example :~sig\\
\hspace*{0pt} ~datatype nat = Zero{\char95}nat | Suc of nat\\
\hspace*{0pt} ~val nat{\char95}aux :~IntInf.int -> nat -> nat\\
\hspace*{0pt} ~val nat :~IntInf.int -> nat\\
\hspace*{0pt} ~type 'a semigroup\\
\hspace*{0pt} ~val mult :~'a semigroup -> 'a -> 'a -> 'a\\
\hspace*{0pt} ~type 'a monoidl\\
\hspace*{0pt} ~val semigroup{\char95}monoidl :~'a monoidl -> 'a semigroup\\
\hspace*{0pt} ~val neutral :~'a monoidl -> 'a\\
\hspace*{0pt} ~type 'a monoid\\
\hspace*{0pt} ~val monoidl{\char95}monoid :~'a monoid -> 'a monoidl\\
\hspace*{0pt} ~type 'a group\\
\hspace*{0pt} ~val monoid{\char95}group :~'a group -> 'a monoid\\
\hspace*{0pt} ~val inverse :~'a group -> 'a -> 'a\\
\hspace*{0pt} ~val pow{\char95}nat :~'a monoid -> nat -> 'a -> 'a\\
\hspace*{0pt} ~val pow{\char95}int :~'a group -> IntInf.int -> 'a -> 'a\\
\hspace*{0pt} ~val mult{\char95}int :~IntInf.int -> IntInf.int -> IntInf.int\\
\hspace*{0pt} ~val semigroup{\char95}int :~IntInf.int semigroup\\
\hspace*{0pt} ~val neutral{\char95}int :~IntInf.int\\
\hspace*{0pt} ~val monoidl{\char95}int :~IntInf.int monoidl\\
\hspace*{0pt} ~val monoid{\char95}int :~IntInf.int monoid\\
\hspace*{0pt} ~val inverse{\char95}int :~IntInf.int -> IntInf.int\\
\hspace*{0pt} ~val group{\char95}int :~IntInf.int group\\
\hspace*{0pt} ~val example :~IntInf.int\\
\hspace*{0pt}end = struct\\
\hspace*{0pt}\\
\hspace*{0pt}datatype nat = Zero{\char95}nat | Suc of nat;\\
\hspace*{0pt}\\
\hspace*{0pt}fun nat{\char95}aux i n =\\
\hspace*{0pt} ~(if IntInf.<= (i,~(0 :~IntInf.int)) then n\\
\hspace*{0pt} ~~~else nat{\char95}aux (IntInf.- (i,~(1 :~IntInf.int))) (Suc n));\\
\hspace*{0pt}\\
\hspace*{0pt}fun nat i = nat{\char95}aux i Zero{\char95}nat;\\
\hspace*{0pt}\\
\hspace*{0pt}type 'a semigroup = {\char123}mult :~'a -> 'a -> 'a{\char125};\\
\hspace*{0pt}val mult = {\char35}mult :~'a semigroup -> 'a -> 'a -> 'a;\\
\hspace*{0pt}\\
\hspace*{0pt}type 'a monoidl = {\char123}semigroup{\char95}monoidl :~'a semigroup,~neutral :~'a{\char125};\\
\hspace*{0pt}val semigroup{\char95}monoidl = {\char35}semigroup{\char95}monoidl :~'a monoidl -> 'a semigroup;\\
\hspace*{0pt}val neutral = {\char35}neutral :~'a monoidl -> 'a;\\
\hspace*{0pt}\\
\hspace*{0pt}type 'a monoid = {\char123}monoidl{\char95}monoid :~'a monoidl{\char125};\\
\hspace*{0pt}val monoidl{\char95}monoid = {\char35}monoidl{\char95}monoid :~'a monoid -> 'a monoidl;\\
\hspace*{0pt}\\
\hspace*{0pt}type 'a group = {\char123}monoid{\char95}group :~'a monoid,~inverse :~'a -> 'a{\char125};\\
\hspace*{0pt}val monoid{\char95}group = {\char35}monoid{\char95}group :~'a group -> 'a monoid;\\
\hspace*{0pt}val inverse = {\char35}inverse :~'a group -> 'a -> 'a;\\
\hspace*{0pt}\\
\hspace*{0pt}fun pow{\char95}nat A{\char95}~Zero{\char95}nat x = neutral (monoidl{\char95}monoid A{\char95})\\
\hspace*{0pt} ~| pow{\char95}nat A{\char95}~(Suc n) x =\\
\hspace*{0pt} ~~~mult ((semigroup{\char95}monoidl o monoidl{\char95}monoid) A{\char95}) x (pow{\char95}nat A{\char95}~n x);\\
\hspace*{0pt}\\
\hspace*{0pt}fun pow{\char95}int A{\char95}~k x =\\
\hspace*{0pt} ~(if IntInf.<= ((0 :~IntInf.int),~k)\\
\hspace*{0pt} ~~~then pow{\char95}nat (monoid{\char95}group A{\char95}) (nat k) x\\
\hspace*{0pt} ~~~else inverse A{\char95}~(pow{\char95}nat (monoid{\char95}group A{\char95}) (nat (IntInf.{\char126}~k)) x));\\
\hspace*{0pt}\\
\hspace*{0pt}fun mult{\char95}int i j = IntInf.+ (i,~j);\\
\hspace*{0pt}\\
\hspace*{0pt}val semigroup{\char95}int = {\char123}mult = mult{\char95}int{\char125}~:~IntInf.int semigroup;\\
\hspace*{0pt}\\
\hspace*{0pt}val neutral{\char95}int :~IntInf.int = (0 :~IntInf.int);\\
\hspace*{0pt}\\
\hspace*{0pt}val monoidl{\char95}int =\\
\hspace*{0pt} ~{\char123}semigroup{\char95}monoidl = semigroup{\char95}int,~neutral = neutral{\char95}int{\char125}~:\\
\hspace*{0pt} ~IntInf.int monoidl;\\
\hspace*{0pt}\\
\hspace*{0pt}val monoid{\char95}int = {\char123}monoidl{\char95}monoid = monoidl{\char95}int{\char125}~:~IntInf.int monoid;\\
\hspace*{0pt}\\
\hspace*{0pt}fun inverse{\char95}int i = IntInf.{\char126}~i;\\
\hspace*{0pt}\\
\hspace*{0pt}val group{\char95}int = {\char123}monoid{\char95}group = monoid{\char95}int,~inverse = inverse{\char95}int{\char125}~:\\
\hspace*{0pt} ~IntInf.int group;\\
\hspace*{0pt}\\
\hspace*{0pt}val example :~IntInf.int =\\
\hspace*{0pt} ~pow{\char95}int group{\char95}int (10 :~IntInf.int) ({\char126}2 :~IntInf.int);\\
\hspace*{0pt}\\
\hspace*{0pt}end;~(*struct Example*)%
\end{isamarkuptext}%
\isamarkuptrue%
%
\endisatagquote
{\isafoldquote}%
%
\isadelimquote
%
\endisadelimquote
%
\begin{isamarkuptext}%
\noindent In Scala, implicts are used as dictionaries:%
\end{isamarkuptext}%
\isamarkuptrue%
%
\isadeliminvisible
%
\endisadeliminvisible
%
\isataginvisible
%
\endisataginvisible
{\isafoldinvisible}%
%
\isadeliminvisible
%
\endisadeliminvisible
%
\isadelimquote
%
\endisadelimquote
%
\isatagquote
%
\begin{isamarkuptext}%
\isatypewriter%
\noindent%
\hspace*{0pt}object Example {\char123}\\
\hspace*{0pt}\\
\hspace*{0pt}abstract sealed class nat\\
\hspace*{0pt}final case object Zero{\char95}nat extends nat\\
\hspace*{0pt}final case class Suc(a:~nat) extends nat\\
\hspace*{0pt}\\
\hspace*{0pt}def nat{\char95}aux(i:~BigInt,~n:~nat):~nat =\\
\hspace*{0pt} ~(if (i <= BigInt(0)) n else nat{\char95}aux(i - BigInt(1),~Suc(n)))\\
\hspace*{0pt}\\
\hspace*{0pt}def nat(i:~BigInt):~nat = nat{\char95}aux(i,~Zero{\char95}nat)\\
\hspace*{0pt}\\
\hspace*{0pt}trait semigroup[A] {\char123}\\
\hspace*{0pt} ~val `Example.mult`:~(A,~A) => A\\
\hspace*{0pt}{\char125}\\
\hspace*{0pt}def mult[A](a:~A,~b:~A)(implicit A:~semigroup[A]):~A =\\
\hspace*{0pt} ~A.`Example.mult`(a,~b)\\
\hspace*{0pt}\\
\hspace*{0pt}trait monoidl[A] extends semigroup[A] {\char123}\\
\hspace*{0pt} ~val `Example.neutral`:~A\\
\hspace*{0pt}{\char125}\\
\hspace*{0pt}def neutral[A](implicit A:~monoidl[A]):~A = A.`Example.neutral`\\
\hspace*{0pt}\\
\hspace*{0pt}trait monoid[A] extends monoidl[A] {\char123}\\
\hspace*{0pt}{\char125}\\
\hspace*{0pt}\\
\hspace*{0pt}trait group[A] extends monoid[A] {\char123}\\
\hspace*{0pt} ~val `Example.inverse`:~A => A\\
\hspace*{0pt}{\char125}\\
\hspace*{0pt}def inverse[A](a:~A)(implicit A:~group[A]):~A = A.`Example.inverse`(a)\\
\hspace*{0pt}\\
\hspace*{0pt}def pow{\char95}nat[A:~monoid](xa0:~nat,~x:~A):~A = (xa0,~x) match {\char123}\\
\hspace*{0pt} ~case (Zero{\char95}nat,~x) => neutral[A]\\
\hspace*{0pt} ~case (Suc(n),~x) => mult[A](x,~pow{\char95}nat[A](n,~x))\\
\hspace*{0pt}{\char125}\\
\hspace*{0pt}\\
\hspace*{0pt}def pow{\char95}int[A:~group](k:~BigInt,~x:~A):~A =\\
\hspace*{0pt} ~(if (BigInt(0) <= k) pow{\char95}nat[A](nat(k),~x)\\
\hspace*{0pt} ~~~else inverse[A](pow{\char95}nat[A](nat((- k)),~x)))\\
\hspace*{0pt}\\
\hspace*{0pt}def mult{\char95}int(i:~BigInt,~j:~BigInt):~BigInt = i + j\\
\hspace*{0pt}\\
\hspace*{0pt}implicit def semigroup{\char95}int:~semigroup[BigInt] = new semigroup[BigInt] {\char123}\\
\hspace*{0pt} ~val `Example.mult` = (a:~BigInt,~b:~BigInt) => mult{\char95}int(a,~b)\\
\hspace*{0pt}{\char125}\\
\hspace*{0pt}\\
\hspace*{0pt}def neutral{\char95}int:~BigInt = BigInt(0)\\
\hspace*{0pt}\\
\hspace*{0pt}implicit def monoidl{\char95}int:~monoidl[BigInt] = new monoidl[BigInt] {\char123}\\
\hspace*{0pt} ~val `Example.neutral` = neutral{\char95}int\\
\hspace*{0pt} ~val `Example.mult` = (a:~BigInt,~b:~BigInt) => mult{\char95}int(a,~b)\\
\hspace*{0pt}{\char125}\\
\hspace*{0pt}\\
\hspace*{0pt}implicit def monoid{\char95}int:~monoid[BigInt] = new monoid[BigInt] {\char123}\\
\hspace*{0pt} ~val `Example.neutral` = neutral{\char95}int\\
\hspace*{0pt} ~val `Example.mult` = (a:~BigInt,~b:~BigInt) => mult{\char95}int(a,~b)\\
\hspace*{0pt}{\char125}\\
\hspace*{0pt}\\
\hspace*{0pt}def inverse{\char95}int(i:~BigInt):~BigInt = (- i)\\
\hspace*{0pt}\\
\hspace*{0pt}implicit def group{\char95}int:~group[BigInt] = new group[BigInt] {\char123}\\
\hspace*{0pt} ~val `Example.inverse` = (a:~BigInt) => inverse{\char95}int(a)\\
\hspace*{0pt} ~val `Example.neutral` = neutral{\char95}int\\
\hspace*{0pt} ~val `Example.mult` = (a:~BigInt,~b:~BigInt) => mult{\char95}int(a,~b)\\
\hspace*{0pt}{\char125}\\
\hspace*{0pt}\\
\hspace*{0pt}def example:~BigInt = pow{\char95}int[BigInt](BigInt(10),~BigInt(- 2))\\
\hspace*{0pt}\\
\hspace*{0pt}{\char125}~/* object Example */%
\end{isamarkuptext}%
\isamarkuptrue%
%
\endisatagquote
{\isafoldquote}%
%
\isadelimquote
%
\endisadelimquote
%
\isamarkupsubsection{Inspecting the type class universe%
}
\isamarkuptrue%
%
\begin{isamarkuptext}%
To facilitate orientation in complex subclass structures, two
  diagnostics commands are provided:

  \begin{description}

    \item[\hyperlink{command.print-classes}{\mbox{\isa{\isacommand{print{\isacharunderscore}classes}}}}] print a list of all classes
      together with associated operations etc.

    \item[\hyperlink{command.class-deps}{\mbox{\isa{\isacommand{class{\isacharunderscore}deps}}}}] visualizes the subclass relation
      between all classes as a Hasse diagram.

  \end{description}%
\end{isamarkuptext}%
\isamarkuptrue%
%
\isadelimtheory
%
\endisadelimtheory
%
\isatagtheory
\isacommand{end}\isamarkupfalse%
%
\endisatagtheory
{\isafoldtheory}%
%
\isadelimtheory
%
\endisadelimtheory
\isanewline
\end{isabellebody}%
%%% Local Variables:
%%% mode: latex
%%% TeX-master: "root"
%%% End:


\begingroup
\bibliographystyle{plain} \small\raggedright\frenchspacing
\bibliography{../../manual}
\endgroup

\end{document}


%%% Local Variables: 
%%% mode: latex
%%% TeX-master: t
%%% End: 
